% Point d'avancement tuteur, vendredi 11 septembre 2026. CINQ pages : titre et quatre slides.
%
% CE QUE CE DECK PORTE, ET IL A ETE VERIFIE PAR GREP SUR TOUS LES DECKS ANTERIEURS. Aucun des
% chiffres de ces quatre slides n'a jamais ete montre : ni le chiffre de tete (6 049, 20 188,
% facteur 3,34), ni la lecture des rapports SFCR (5 305, 812, 16,7 %), ni la dynamique du
% facteur systemique. Le dernier point PRESENTE est celui du 14 aout ; le deck du 21 aout porte
% d'autres sujets (defaillances simultanees, additivites, concavite en duree).
%
% LE FIL : deux semaines d'absence, donc le deck ne raconte pas une production de semaine. Il
% porte le chiffre de tete, qui n'avait jamais ete expose, et les deux verifications de fond
% faites le 2 septembre. C'est un deck de CONSOLIDATION, et il faut l'annoncer comme tel.
%
% HORS SUJET SUR CONSIGNE, NE PAS REMETTRE : la charte graphique, et tout ce qui releve de la
% FABRICATION du document (comptes de pages, taux du harnais, defauts de citation, couverture
% de deck). Ce travail est trace dans CLAUDE.md et PLAN_MEMOIRE.md.
%
% REGLE DE DENSITE : support de parole, pas document. Titre qui enonce le resultat, une table ou
% trois a cinq lignes courtes, les chiffres. Le raisonnement passe a l'oral. Les tables valent
% mieux que les puces des que la slide compare deux etats. NE PAS REDENSIFIER.
%
% AUCUNE FIGURE. La semaine a produit des verifications et une section de modele, pas de
% nouvelles courbes. S10 (SCR par etat) et S24 (interaction des canaux) restent disponibles,
% aucune n'ayant jamais servi dans un deck.
%
% PUCES : template beamer force en cercle. Le babel francais met un tiret cadratin en puce
% d'itemize par defaut, et les tirets cadratins sont proscrits. Controle SUR LE PDF.
%
% HARNAIS : chaque section porte sa ligne de sources, la page de titre est declaree hors
% section. Une declaration n'est honoree que si la section ne cite aucun script, et elle doit
% etre posee APRES le \section*.
\documentclass[11pt]{beamer}
\usetheme{Madrid}
\usecolortheme{seahorse}
\usepackage[utf8]{inputenc}
\usepackage[T1]{fontenc}
\usepackage[french]{babel}
\usepackage{booktabs}
\usepackage{amsmath,amssymb}
\usepackage{eurosym}
\graphicspath{{../vasicek_lab/figures/}{../cascade_qualitative/figures/}{../vasicek_lab/notes/}}
\setbeamertemplate{navigation symbols}{}
\setbeamertemplate{caption}{\raggedright\insertcaption\par}
\setbeamertemplate{itemize items}[circle]
\setbeamerfont{block title}{size=\small}

\title[Point d'avancement]{Mémoire SCR DORA : point d'avancement}
\subtitle{Le chiffre de tête, et deux vérifications de fond}
\author[K. Kaddouri]{Kélian Kaddouri}
\institute[ENSAE / Nexialog]{ENSAE / Nexialog Consulting}
\date[11/09/2026]{11 septembre 2026}

\begin{document}

% HARNAIS-HORS-SECTION: page de titre. Millesime et date, aucun resultat.
\begin{frame}
\titlepage
\end{frame}

% =====================================================================
\section*{A. Le chiffre de tête}

\begin{frame}{Quatre protocoles chiffraient l'écart. Un seul est publié}
\small
Quatre protocoles du mémoire chiffrent \og{}l'écart entre l'état conforme et l'état non
conforme\fg{}. Aucun n'est faux : ils ne relâchent pas le même nombre de canaux.
\vspace{3mm}

\begin{center}\small
\begin{tabular}{@{}l c r r r@{}}
\toprule
\textbf{Lecture} & \textbf{canaux} & \textbf{conforme} & \textbf{non conf.} & \textbf{facteur} \\
\midrule
Score de conformité seul & $1$ & $7\,987$ & $16\,847$ & $2{,}11$ \\
Fréquence $+$ propagation & $2$ & $6\,664$ & $15\,074$ & $2{,}26$ \\
$+$ détection & $3$ & $5\,932$ & $16\,595$ & $2{,}80$ \\
\textbf{$+$ accumulation tiers} & $\mathbf{4}$ & $\mathbf{6\,049}$ & $\mathbf{20\,188}$ &
$\mathbf{3{,}34}$ \\
\bottomrule
\end{tabular}
\end{center}
\vspace{4mm}

\begin{itemize}\itemsep6pt
  \item \alert{La quatrième est la seule où l'état conforme est conforme sur tous les canaux}
  \item Ailleurs, c'est une remédiation partielle appelée conformité, ce qui sous-estime
        l'écart par construction
  \item Les trois autres restent publiées, comme remédiations partielles
\end{itemize}
\vspace{2mm}
{\scriptsize En M\euro, périmètre OpRisk. Sources : scripts~\texttt{16}, \texttt{25},
\texttt{43} et~\texttt{67}.}
\end{frame}

% =====================================================================
\section*{B. Les quatre bilans réels}

\begin{frame}{Les rapports SFCR relus, et une erreur de champ}
\small
Les quatre jeux de chiffres d'entité étaient saisis depuis les rapports de solvabilité 2024
sans que les rapports eux-mêmes aient été relus. Ils l'ont été.
\vspace{3mm}

\begin{center}\small
\begin{tabular}{@{}l r r@{}}
\toprule
\textbf{Assureur non-vie B} & \textbf{avant} & \textbf{après relecture} \\
\midrule
Total d'actif & $3\,234$ & $\mathbf{5\,305}$ \\
\textsc{scr} & $292$ \emph{(déduit)} & $\mathbf{812{,}5}$ \emph{(publié)} \\
Besoin \textsc{orsa} au titre de DORA & $120{,}8$ & $135{,}5$ \\
Rapport au \textsc{scr} publié & $41{,}3\,\%$ & $\mathbf{16{,}7\,\%}$ \\
\bottomrule
\end{tabular}
\end{center}
\vspace{3mm}

\begin{itemize}\itemsep5pt
  \item Le $3\,234$ était saisi en \textbf{total d'actif} là où le rapport le donne en
        \textbf{fonds propres éligibles}
  \item L'identité comptable le confirme : $5\,305 - 2\,071 = 3\,234$
  \item \alert{La figure signalait l'anomalie avant qu'on la trouve} : la part montait avec la
        taille entre les deux petites entités
\end{itemize}
\vspace{1mm}
{\scriptsize En M\euro. Sources : script~\texttt{65}.}
\end{frame}

% =====================================================================
\begin{frame}{Ce que cette relecture fait gagner}
\small
\begin{block}{Une limite déclarée disparaît}
Les quatre \textsc{scr} sont \textbf{publiés} par les rapports. Plus aucun n'est reconstitué
depuis un taux de couverture, donc la réserve \og{}deux \textsc{scr} sur quatre sont
déduits\fg{} et la borne sous stress de $\pm 10\,\%$ n'ont plus d'objet.
\end{block}
\vspace{3mm}
\begin{itemize}\itemsep6pt
  \item \textbf{Le forfait de $3\,\%$ n'est pas une borne, et c'est mesuré} : contre le module
        opérationnel publié, $54{,}1$ contre $59$ chez l'un, $62{,}1$ contre $39{,}7$ chez
        l'autre. Il se trompe \alert{dans les deux sens}
  \item \textbf{Le test du levier a changé de sens}, et je ne l'invoque plus : la part de
        l'entité notionnelle passe de $13$ à $9\,\%$, donc \emph{sous} le critère de $10\,\%$
  \item La requalification du $169$~M\euro{} repose désormais sur le \textbf{seul} argument de
        taille, qui tient par lui-même
\end{itemize}
\vspace{2mm}
{\scriptsize Sources : script~\texttt{65}.}
\end{frame}

% =====================================================================
\section*{C. La dynamique du facteur}

\begin{frame}{Le facteur systémique suit-il une trajectoire ?}
\small
Question de ma tutrice de stage : sous Vasicek, la probabilité de défaut de crédit est portée
par un brownien. Quelle trajectoire suit mon facteur systémique ?
\vspace{3mm}

\begin{itemize}\itemsep6pt
  \item \textbf{Réponse en deux temps.} Le facteur publié n'a \emph{pas} de trajectoire : il est
        indépendant d'un pas au suivant. Ce n'est pas un défaut, c'est \alert{exactement
        l'incrément brownien normalisé} du modèle de crédit, au pas comptable
  \item La dépendance entre piliers existe et elle est forte : charge commune $0{,}68$, soit une
        corrélation de $0{,}462$ entre deux piliers
  \item \textbf{Mais un brownien pur imposerait un indice de dispersion égal à $1$}, et il est
        mesuré à $1{,}19$ sur les cellules firme-année et $2{,}24$ sur la série annuelle. Le
        moteur en pose $9{,}20$ comme marge
  \item \alert{Donc la donnée identifie une composante de saut} : le facteur s'écrit avec un
        retour à la moyenne et des sauts, et la section le publie avec sa normalisation
\end{itemize}
\vspace{3mm}
{\scriptsize L'indice de dispersion vaut $1 + \lambda/r$ : il n'est pas invariant d'échelle,
d'où trois valeurs pour trois niveaux d'agrégation. Sources : scripts~\texttt{08b},
\texttt{16}, \texttt{41} et~\texttt{69}.}
\end{frame}

\end{document}
